\section{Що агрегований інтерфейс зберігає, а що втрачає}
\label{sec:analysis}

\subsection{Представлення за допомогою сурогатної мережі}
\label{sec:surrogate}

За точного обчислення часу~\eqref{eq:hazard_integral} ідеальний двомасштабний процес є CTMC на $\{S,I,R\}^V$. Інтенсивність його переходу $x\to x^{v:I}$ унаслідок імпорту дорівнює $\sum_{i\neq j}\lambda_{ij}(x)\Pi_{ij}(v\mid x)$ для $v\in S_j(x)$; усі локальні переходи мають інтенсивності з~\eqref{eq:full_generator}, обмежені на $E_{ii}$. Наступне твердження ототожнює цей процес із процесом на повній мережі зі зміненим графом.

\begin{proposition}[Сурогатна мережа]
\label{prop:surrogate}
Нехай $\tilde G$ --- мережа з тими самими внутрішньоспільнотними каналами та інтенсивностями, що й $G$, у якій кожен міжспільнотний пучок $E_{ij}$ замінено повною множиною каналів $V_i\times V_j$, кожен з інтенсивністю $\tilde\beta_{uv}=\beta_{\mathrm{out}}W_{ij}/(N_iN_j)$. За рівномірного розміщення ідеальний двомасштабний процес має генератор $\mathcal{L}_{\tilde G}$; тобто він має закон процесу на повній мережі $\tilde G$, запущеного з тієї самої конфігурації.
\end{proposition}

\begin{proof}
Видалення та внутрішньоспільнотні зараження мають однакові інтенсивності в обох процесах. Розглянемо сприйнятливий вузол $v\in V_j$ у конфігурації $x$. За рівномірного розміщення двомасштабна інтенсивність, з якою $v$ заражається внаслідок імпорту, дорівнює
\[
\sum_{i\neq j}\lambda_{ij}(x)\,\frac{1}{s_j(x)}
=\sum_{i\neq j}\beta_{\mathrm{out}}W_{ij}\frac{y_i(x)}{N_i}\frac{s_j(x)}{N_j}\frac{1}{s_j(x)}
=\sum_{i\neq j}\sum_{u\in V_i}\frac{\beta_{\mathrm{out}}W_{ij}}{N_iN_j}\ind\{x_u=I\},
\]
що є сумарною інтенсивністю міжспільнотних каналів $\tilde G$, спрямованих у $v$. Коли $s_j(x)=0$, обидві інтенсивності дорівнюють нулю. Отже, всі інтенсивності переходів збігаються, а генератори є однаковими. \qed
\end{proof}

\begin{corollary}
\label{cor:exact}
Якщо кожен міжспільнотний пучок у $G$ уже є повним, $E_{ij}=V_i\times V_j$, зі спільною інтенсивністю для кожної пари спільнот, тоді $\tilde G=G$, а ідеальний двомасштабний процес є точним.
\end{corollary}

Сурогатна мережа зберігає сумарну вагу кожного пучка, $\sum_{(u,v)\in V_i\times V_j}\tilde\beta_{uv}=\beta_{\mathrm{out}}W_{ij}$, але рівномірно розподіляє її між усіма парами (рисунок~\ref{fig:schematic}c). Тому похибка апроксимації ідеального методу є різницею між епідеміями на $G$ та $\tilde G$, які відрізняються лише своїми інтерфейсами. Реалізований метод додає чисельну похибку синхронізації. Твердження~\ref{prop:surrogate} також дає практичний спосіб верифікації: для малих мереж похибку синхронізації можна виокремити, порівнюючи двомасштабний симулятор із точною симуляцією на $\tilde G$.

\subsection{Джерела розбіжностей}
\label{sec:sources}

Порівняння $G$ із $\tilde G$ виявляє чотири механізми, через які двомасштабний симулятор може відхилятися від процесу на повній мережі.

\paragraph{Інформація про стан межі (з боку джерела).}
У $G$ інтенсивність передавання з $V_i$ до $V_j$ дорівнює
\begin{equation}
\label{eq:exact_rate}
a^\star_{ij}(x)=\beta_{\mathrm{out}}\sum_{(u,v)\in E_{ij}}w_{uv}\,\ind\{x_u=I,\,x_v=S\},
\end{equation}
і залежить від того, чи заражені конкретні мостові вузли. У $\tilde G$ свій внесок робить кожен заражений вузол $V_i$. Тому експорти можуть відбуватися до зараження будь-якого мостового вузла. Такий \emph{обхід мостів} неможливий у $G$ і дає безпосередній діагностичний показник (розділ~\ref{sec:single_export}).

\paragraph{Розміщення (з боку приймача).}
У $G$ передавання заражає цільову кінцеву вершину каналу, що спрацював; у $\tilde G$ ціль рівномірно вибирається серед сприйнятливих вузлів. Місце імпорту визначає його відстань до наступного інтерфейсу, а отже, і час подальших переходів.

\paragraph{Вичерпання інтерфейсу (з боку приймача).}
За SI-динаміки канал $G$ передає зараження не більш ніж один раз: після цього його ціль заражена. Тому пучок із $|E_{ij}|$ каналів переносить не більш ніж $|E_{ij}|$ заражень, а інтерфейс з одним ребром --- рівно одне. У $\tilde G$ імпорти з $V_i$ у $V_j$ тривають з інтенсивністю $\lambda_{ij}=\beta_{\mathrm{out}}W_{ij}\ibar_i\sbar_j$, доки $V_j$ не буде вичерпано; їхня очікувана кількість за період $[t_1,t_2]$ дорівнює $\mathbb{E}\int_{t_1}^{t_2}\lambda_{ij}(s)\,ds$. Отже, коли приймальна спільнота тривалий час залишається переважно сприйнятливою, двомасштабний симулятор продовжує засівати її додатковими зараженнями, яких повна мережа передати не може.

\paragraph{Синхронізація.}
Квадратурна формула середньої точки та інтерполяція з розділу~\ref{sec:timing} апроксимують~\eqref{eq:hazard_integral}.

\begin{remark}[Розміщення за рандомізованих інтерфейсів]
\label{rem:placement}
Припустімо, що кінцеві вершини кожного пучка вибираються рівномірно й незалежно від внутрішніх графів та кінцевих вершин інших пучків, як у наведених нижче експериментах. Тоді ціль першого передавання до повністю сприйнятливої спільноти має рівномірний розподіл і в $G$, оскільки вона є кінцевою вершиною рівномірно розміщеного каналу, і в $\tilde G$. Тому розміщення може впливати на часи прибуття лише через передавання після першого. У таких схемах відмінності між першим і подальшими переходами виокремлюють ефекти з боку приймача (розділ~\ref{sec:res_descriptor}).
\end{remark}

\subsection{Одиничний експорт: стохастичне впорядкування та обхід мостів}
\label{sec:single_export}

Розглянемо спільноту-джерело $V_0$ розміру $n$ із довільним внутрішнім графом, SI-динамікою та єдиним сприйнятливим листковим вузлом $\ell$, сполученим із множиною $B\subset V_0$ з $k$ мостових вузлів каналами з інтенсивністю $\beta_{\mathrm{out}}$. Тоді $W_{0\ell}=k$, а до експорту $\sbar_\ell=1$, тому $\lambda_{0\ell}=\beta_{\mathrm{out}}k\,\ibar_0$. Нехай $\tau_u$ --- час зараження вузла $u$, $A_u(t)=(t-\tau_u)^+$ --- тривалість його перебування в зараженому стані, а $T_B=\min_{u\in B}\tau_u$.

\begin{proposition}[Агреговані годинники експорту спрацьовують раніше]
\label{prop:ordering}
Нехай $B$ має рівномірний розподіл на множині всіх $k$-елементних підмножин $V_0$ незалежно від внутрішнього графа, початкової конфігурації та випадковості подій. Нехай $T^{\Full}$ і $T^{\TS}$ --- часи першого зараження $\ell$ відповідно в повній мережі та в ідеальному двомасштабному процесі. Тоді
\[
\mathbb{P}(T^{\Full}>t)\;\ge\;\mathbb{P}(T^{\TS}>t)\qquad\text{for all }t\ge0,
\]
отже, $T^{\TS}$ є стохастично меншим за $T^{\Full}$ \citep{shaked2007}. Зокрема, середнє та кожен квантиль часу експорту ідеального двомасштабного процесу не перевищують відповідних величин часу експорту в повній мережі.
\end{proposition}

\begin{proof}
До зараження $\ell$ не впливає на $V_0$: він нічого не передає, а інтенсивності вузлів джерела залежать лише від їхніх заражених сусідів. Тому обидва часи експорту можна побудувати на одному ймовірнісному просторі з автономної SI-траєкторії $X_0$ спільноти-джерела, продовженої після експорту, та незалежного порога $E\sim\Exp(1)$:
\[
T^{\Full}=\inf\Bigl\{t:\beta_{\mathrm{out}}\sum_{u\in B}A_u(t)\ge E\Bigr\},\qquad
T^{\TS}=\inf\Bigl\{t:\beta_{\mathrm{out}}\frac{k}{n}\sum_{u\in V_0}A_u(t)\ge E\Bigr\},
\]
де другий підінтегральний вираз випливає з $\int_0^t\ibar_0(s)\,ds=n^{-1}\sum_uA_u(t)$. Умовні щодо $X_0$ ймовірності виживання дорівнюють $\exp(-\beta_{\mathrm{out}}\sum_{u\in B}A_u(t))$ та $\exp(-\beta_{\mathrm{out}}(k/n)\sum_uA_u(t))$. Оскільки $B$ є рівномірною та незалежною від $X_0$, маємо $\mathbb{E}[\sum_{u\in B}A_u(t)\mid X_0]=(k/n)\sum_uA_u(t)$, а нерівність Єнсена для опуклої функції $z\mapsto e^{-\beta_{\mathrm{out}}z}$ дає $\mathbb{P}(T^{\Full}>t\mid X_0)\ge\mathbb{P}(T^{\TS}>t\mid X_0)$. Взяття математичного сподівання завершує доведення. \qed
\end{proof}

Нерівність є строгою, якщо з додатною ймовірністю тривалості перебування вузлів у зараженому стані в момент $t$ відрізняються та $0<k<n$. Вона також виконується, коли $B$ рівномірно розподілена серед вузлів, відмінних від початково зараженого, оскільки вилучення вузла, зараженого найдавніше, може лише зменшити умовне середнє мостової небезпеки. Величина розриву визначається умовною дисперсією мостової небезпеки, для якої вибірка без повернення дає
\begin{equation}
\label{eq:hazard_variance}
\operatorname{Var}\Bigl(\beta_{\mathrm{out}}\sum_{u\in B}A_u(t)\Bigm| X_0\Bigr)=\beta_{\mathrm{out}}^2\,k\,\frac{n-k}{n-1}\,\operatorname{Var}_u\bigl(A_u(t)\bigr),
\end{equation}
де $\operatorname{Var}_u$ --- генеральна дисперсія за вузлами джерела. Після повного зараження джерела $\operatorname{Var}_u(A_u(t))=\operatorname{Var}_u(\tau_u)$. Отже, агрегований годинник є близьким до точного, коли часи зараження вузлів зосереджені порівняно з $1/\beta_{\mathrm{out}}$, як у кліці, що швидко насичується, і спрацьовує раніше, коли зараження поширюється джерелом повільно й нерівномірно.

Та сама побудова дає ймовірність обходу мостів. У $G$ експорт не може передувати $T_B$. В ідеальному двомасштабному процесі
\begin{equation}
\label{eq:bypass}
\mathbb{P}\bigl(T^{\TS}<T_B\bigr)=\mathbb{E}\Bigl[1-\exp\Bigl(-\beta_{\mathrm{out}}\frac{k}{n}\int_0^{T_B}y_0(s)\,ds\Bigr)\Bigr],
\end{equation}
що є ймовірністю того, що агрегована небезпека, накопичена до активації справжнього інтерфейсу, спричинить експорт. Вона зростає, коли зараження досягає мостів пізно та після багатьох інших вузлів, як це відбувається в розріджених спільнотах із довгими шляхами. Додаток~\ref{app:clique} наводить $\mathbb{E}[T_B]$ для джерела-кліки.

\subsection{Характеристика співвідношення часових масштабів}
\label{sec:descriptor}

Механізми як з боку джерела, так і з боку приймача діють лише доти, доки зараження ще поширюється всередині спільноти. Згідно з~\eqref{eq:hazard_variance}, розрив з боку джерела визначається добутком $\beta_{\mathrm{out}}$ і розсіювання часів зараження. Очікувана кількість надлишкових імпортів до спільноти під час внутрішнього поширення має порядок $\beta_{\mathrm{out}}W_{ij}$, помножений на тривалість цього поширення. Нехай $\sigma_T$ позначає стандартне відхилення часів зараження вузлів усередині спільноти, яке вимірює цю тривалість. Як характеристику очікуваної точності ми використовуємо безрозмірне співвідношення
\begin{equation}
\label{eq:rho}
\rho=\beta_{\mathrm{out}}\,\sigma_T,
\end{equation}
тобто внутрішній розкид часу в одиницях середнього часу переходу через одне ребро $1/\beta_{\mathrm{out}}$. Ці два механізми по-різному масштабуються із шириною інтерфейсу ($\sqrt{k}$ у~\eqref{eq:hazard_variance} та $k$ для надлишкових імпортів, що частково компенсується тим, що сама повна мережа допускає до $k$ переходів). Тому ми не включаємо $k$ до $\rho$ та перевіряємо цю характеристику емпірично, а не виводимо закон похибки.

%======================================================================
\section{Дизайн експерименту}
\label{sec:design}

\subsection{Тестові задачі}

У таблиці~\ref{tab:design} підсумовано чотири серії експериментів. Усі вони використовують SI-динаміку, одиничні ваги контактів, однакові внутрішні й зовнішні інтенсивності на ребро ($\beta_{\mathrm{in}}=\beta_{\mathrm{out}}=\beta$) та планувальник з алгоритму~\ref{alg:scheduler} з $\tau=1$. У тестовій задачі \emph{кліка--листок} листок приєднано до $k=2$ мостових вузлів кліки зі 100 вузлів. У \emph{переборі джерело--листок} кліку замінено зв'язними випадковими графами з $n=1000$ вузлів із заданою середньою довжиною найкоротшого шляху (ASPL). У серії \emph{ланцюг} $L=10$ спільнот з'єднано в лінію за допомогою $k$ ребер між послідовними спільнотами; зараження починається у спільноті 0. Зв'язні випадкові графи генеруються додаванням рівномірно випадкових ребер до рівномірно випадкового дерева з міченими вершинами; кількість доданих ребер калібрується так, щоб ASPL лежала в межах 0.1 від цільового значення; ASPL 1 відповідає повному графу. Тому середній степінь зменшується зі зростанням цільової ASPL. За цільової ASPL 6 він змінюється від 2.08 для $n=50$ до 3.4 для $n=1000$. Кінцеві вершини інтерфейсів вибираються рівномірно й незалежно для кожного інтерфейсу та повтору.\footnote{У переборі джерело--листок множини мостів вибираються з усіх вузлів, крім одного вузла з фіксованою міткою. Оскільки випадкові графи та рівномірно вибране початкове зараження є інваріантними за законом відносно перемічування міток, це еквівалентно рівномірному розміщенню, тож припущення твердження~\ref{prop:ordering} виконується.} Початковий заражений вузол є псевдовипадковим і спільним для обох симуляторів; він може бути мостовим вузлом. У тестовій задачі кліка--листок початковим є фіксований вузол кліки.

\begin{table}[tbp]
\centering
\caption{Серії експериментів. Кожна комірка перебору використовує 40 реалізацій мережі (блоків) із 10 повторами інтерфейсу та подій у кожній. ASPL: цільова внутрішня середня довжина найкоротшого шляху (1 позначає повний граф).}
\label{tab:design}
\small
\begin{tabular}{>{\raggedright\arraybackslash}p{0.17\linewidth}>{\raggedright\arraybackslash}p{0.36\linewidth}>{\raggedright\arraybackslash}p{0.06\linewidth}>{\raggedright\arraybackslash}p{0.08\linewidth}>{\raggedright\arraybackslash}p{0.2\linewidth}}
\toprule
Серія & Структура & $k$ & $\beta$ & Запусків на комірку та модель\\
\midrule
Кліка--листок & кліка-джерело $n=100$ $+$ один листок & 2 & 0.005 & 10\,000\\
Перебір джерело--листок & джерело $n=1000$, ASPL 1--6 із кроком 0.25, $+$ один листок & 1, 5 & 0.003 & 400 (40 $\times$ 10)\\
Ланцюг клік & $L=10$ клік із $n=1000$ & 1 & 0.003 & 400 (40 $\times$ 10)\\
Перебір розміру ланцюга & $L=10$, $n\in\{50,100,250,1000\}$, ASPL 1--6 із кроком 1 & 1, 5 & 0.003 & 400 (40 $\times$ 10)\\
\bottomrule
\end{tabular}
\end{table}

Запуски \Full{} та \TS{} є зіставленими: вони мають спільні граф, інтерфейс і початковий вузол, але використовують незалежні потоки подій. Усі експорти та прибуття, що входять до метрик, було спостережено, тому поправка на цензурування не потрібна. Горизонт становив 1000 для тестової задачі кліка--листок та 30\,000 для решти. Два запуски \Full{} було продовжено до 40\,000, щоб спостерегти їхні термінальні прибуття (додаток~\ref{app:repro}). Ланцюг клік є коміркою ASPL-1, $n=1000$, $k=1$ перебору розміру, поданою окремо, оскільки це найсприятливіший випадок.

\subsection{Спостережувані величини та метрики}

Для тестових задач із листком ми реєструємо час експорту $T_{\mathrm{export}}$ (перше зараження листка) та час активації мосту $T_B$. Для ланцюгів $T_c$ --- час першого зараження спільноти $c$, відрахований від початкового зараження; термінальне прибуття --- це $T_9$. Ми подаємо три підсумкові показники:
\begin{itemize}
\item відносний зсув медіани $b_c=100\,[\mathrm{med}(T^{\TS}_c)-\mathrm{med}(T^{\Full}_c)]/\mathrm{med}(T^{\Full}_c)$, де від'ємні значення означають раніше прибуття в \TS{};
\item відносну відстань Вассерштайна-1 $\eW=100\,W_1/\overline{T}^{\Full}$, де $W_1=R^{-1}\sum_r|T^{\TS}_{(r)}-T^{\Full}_{(r)}|$ обчислюється за порядковими статистиками $R$ часів кожної моделі;
\item частку обходу мостів, тобто відсоток переходів у \TS{}, що відбуваються до зараження будь-якого мостового вузла спільноти-джерела.
\end{itemize}
Для ланцюгів \emph{локальний} зсув переходу $c$ порівнює середні значення $T_c-T_{c-1}$ між моделями. Ми зіставляємо перший перехід з об'єднаними переходами 2--9, використовуючи $T_9-T_1$.

\subsection{Статистичні процедури}

Довірчі інтервали є 95\% процентильними інтервалами з бутстрепу, який повторно вибирає цілі мережеві блоки, зберігаючи разом їхні десять повторів; для підсумкових показників ланцюга використовуються 2000 бутстреп-вибірок. Оскільки $\eW$ є додатною навіть тоді, коли дві вибірки походять з одного закону розподілу, ми калібруємо її за допомогою перестановочного тесту всередині пар: мітки двох запусків у кожній зіставленій парі міняються місцями з імовірністю $1/2$, після чого $\eW$ повторно обчислюється 2000 разів. За нульової гіпотези про обмінність зіставлених запусків це дає вибірковий нижній поріг, який ми підсумовуємо його 95-м процентилем, а також $p$-значення для спостережуваної $\eW$. Додаток~\ref{app:floor} наводить аналітичне наближення цього порога, застосовне, коли дані окремих запусків недоступні.

Початкові значення генератора подій для повтору залежать від цільової ASPL, мережевого блоку та індексу повтору, але не від $n$ або $k$; отже, початкові значення є спільними для різних розмірів спільноти та ширин інтерфейсу (спільні випадкові числа). Це зменшує дисперсію порівнянь між такими комірками, але також корелює їхні похибки вибірки. Тому ми не вважаємо узгодженість між комірками незалежним повторенням.

\subsection{Верифікація еталонної моделі}

Для джерела-кліки з рівномірно розміщеними мостами очікуваний час активації мосту відомий у замкненій формі (додаток~\ref{app:clique}). Для чотирьох розмірів ланцюга клік і обох ширин інтерфейсу середнє $T_B$ у \Full{} узгоджувалося з формулою в межах 2.3 стандартної похибки. Для $k=1$ середній час очікування після активації мосту (326.5--354.0) узгоджувався з точним значенням $1/\beta=333.3$ у межах похибки вибірки (таблиця~\ref{tab:verification}).

%======================================================================
